\documentclass[a4paper,11pt]{article}
\usepackage[utf8]{inputenc}
\usepackage[T1]{fontenc}
\usepackage{lmodern}
\usepackage{amsmath,amssymb,amsthm}
\usepackage{mathtools}
\usepackage{booktabs}
\usepackage{longtable}
\usepackage{array}
\usepackage{hyperref}
\usepackage{graphicx}
\usepackage{float}
\usepackage{geometry}
\geometry{margin=25mm}

\hypersetup{
  colorlinks=true,
  linkcolor=blue,
  urlcolor=blue,
  citecolor=blue
}

\theoremstyle{definition}
\newtheorem{theorem}{Theorem}[section]
\newtheorem{proposition}[theorem]{Proposition}
\newtheorem{lemma}[theorem]{Lemma}
\newtheorem{corollary}[theorem]{Corollary}
\newtheorem{definition}[theorem]{Definition}
\newtheorem{remark}[theorem]{Remark}
\newtheorem{observation}[theorem]{Observation}

\setlength{\parindent}{1em}
\setlength{\parskip}{0.5em}

\providecommand{\tightlist}{\setlength{\itemsep}{0pt}\setlength{\parskip}{0pt}}
\begin{document}

\section{Schwarzschild--de Sitter Exact Solution in the Central
Projection
Framework}\label{schwarzschildde-sitter-exact-solution-in-the-central-projection-framework}

\textbf{--- A Spherically Symmetric Static Exact Solution of the Vacuum
Einstein Equation from Paper {[}1{]} ---}

\begin{center}\rule{0.5\linewidth}{0.5pt}\end{center}

\textbf{Author}: Noriaki Kihara (WF System Co., Ltd.)

\textbf{Date}: April 2026

\textbf{DOI:}
\href{https://doi.org/10.5281/zenodo.19538098}{10.5281/zenodo.19538098}

\begin{center}\rule{0.5\linewidth}{0.5pt}\end{center}

\subsection{Abstract}\label{abstract}

In this paper, we construct the Schwarzschild--de Sitter metric as the
spherically symmetric static exact solution of the vacuum Einstein
equation \(G_{\mu\nu} + \Lambda\, g_{\mu\nu} = 0\) (\(\Lambda = 3/R^2\))
derived in the preceding paper {[}1{]}. This exact solution is uniquely
determined from the equation in {[}1{]} without additional assumptions
(by Birkhoff's theorem). We take the limit \(R \to 0\) in the
constructed exact solution and describe how the structural tension
identified in paper {[}4{]} is expressed in the language of this exact
solution.

\textbf{Disclaimer}: This paper only constructs a mathematically
rigorous solution of the equation derived in paper {[}1{]} and does not
claim that this exact solution describes any actual astrophysical
object. The Schwarzschild--de Sitter metric is a well-known solution in
general relativity, and this paper positions that solution within the
context of the central projection framework. Physical interpretation is
outside the scope of this paper and is left to the reader.

\begin{center}\rule{0.5\linewidth}{0.5pt}\end{center}

\subsection{1. Introduction}\label{introduction}

In the preceding paper {[}1{]}, Proposition 6.1, the vacuum Einstein
equation

\[G_{\mu\nu} + \Lambda\, g_{\mu\nu} = 0, \qquad \Lambda = \frac{3}{R^2} \tag{1.1}\]

was derived for the induced metric \(g_{\mu\nu}\) under the central
projection \(\Phi : \Pi_R \to S^4(R)\). This equation was proved as a
geometric identity (paper {[}1{]}, Remark 6.1) and contains no
additional physical assumptions.

The purpose of this paper is to construct the spherically symmetric
static exact solution of equation (1.1) and to examine the parameter
dependence on the curvature radius \(R\), particularly the behavior in
the limit \(R \to 0\).

\begin{center}\rule{0.5\linewidth}{0.5pt}\end{center}

\subsection{2. Verification of
Prerequisites}\label{verification-of-prerequisites}

Before seeking the spherically symmetric static exact solution of
equation (1.1), we verify the structural properties of the central
projection framework.

\subsubsection{2.1 Spherical Symmetry}\label{spherical-symmetry}

The central projection \(\Phi\) is defined as a projection onto the
hypersphere \(S^n(R)\) of radius \(R\) in the background space
\(\mathbb{R}^{n+1}\) (paper {[}1{]}, §2). This construction is invariant
under rotations about the central axis and possesses spherical symmetry.

\subsubsection{2.2 Vacuum Condition}\label{vacuum-condition}

The right-hand side of equation (1.1) is zero. That is, it is a vacuum
equation that does not include contributions from the energy-momentum
tensor \(T_{\mu\nu}\).

\subsubsection{\texorpdfstring{2.3 Origin of the Cosmological Term
\(\Lambda\)}{2.3 Origin of the Cosmological Term \textbackslash Lambda}}\label{origin-of-the-cosmological-term-lambda}

\(\Lambda = 3/R^2\) is a quantity derived from the geometric structure
of the central projection (paper {[}1{]}, equation 6.2). Unlike the
cosmological constant introduced externally in general relativity, in
the present framework \(\Lambda\) is \textbf{intrinsically determined}
as a function of the curvature radius \(R\).

\subsubsection{2.4 Stationarity}\label{stationarity}

In the construction of paper {[}1{]}, the subjective space \(\Pi_R\)
contains neither positional nor temporal information. The induced metric
\(g_{\mu\nu}\) depends only on the coordinates \(x^\mu\) and does not
contain a time variable. This property is consistent with the exact
solution being static.

\begin{center}\rule{0.5\linewidth}{0.5pt}\end{center}

\subsection{3. Construction of the Schwarzschild--de Sitter Exact
Solution}\label{construction-of-the-schwarzschildde-sitter-exact-solution}

\subsubsection{3.1 Application of Birkhoff's
Theorem}\label{application-of-birkhoffs-theorem}

By the \(\Lambda\)-extended version of Birkhoff's theorem in general
relativity (Kottler, 1918), the solution of the spherically symmetric
vacuum Einstein equation

\[G_{\mu\nu} + \Lambda\, g_{\mu\nu} = 0 \tag{3.1}\]

is static and uniquely reduces to the following Schwarzschild--de Sitter
metric (Kottler metric):

\[ds^2 = -f(r)\,dt^2 + f(r)^{-1}\,dr^2 + r^2\,d\Omega^2 \tag{3.2}\]

where

\[f(r) = 1 - \frac{2M}{r} - \frac{\Lambda}{3}\,r^2 \tag{3.3}\]

Here \(M\) is an integration constant (mass parameter) and \(d\Omega^2\)
is the metric on the unit sphere.

\subsubsection{3.2 Substitution into the Central Projection
Framework}\label{substitution-into-the-central-projection-framework}

Substituting \(\Lambda = 3/R^2\) from equation (1.1) into equation (3.3)
yields

\[f(r) = 1 - \frac{2M}{r} - \frac{r^2}{R^2} \tag{3.4}\]

This is the Schwarzschild--de Sitter exact solution in the central
projection framework.

\subsubsection{3.3 Special Cases}\label{special-cases}

\textbf{Case \(M = 0\)}: Equation (3.4) becomes

\[f(r) = 1 - \frac{r^2}{R^2} \tag{3.5}\]

reducing to the pure de Sitter metric. This corresponds to the case
without a mass source, and represents the spherically symmetric
expression of the induced metric of the central projection itself.

\textbf{Case \(R \to \infty\) (\(\Lambda \to 0\))}: Equation (3.4)
becomes

\[f(r) \to 1 - \frac{2M}{r} \tag{3.6}\]

reducing to the standard Schwarzschild metric. This is consistent with
the \(R \to \infty\) limit in paper {[}4{]}, §2.

\begin{center}\rule{0.5\linewidth}{0.5pt}\end{center}

\subsection{\texorpdfstring{4. Examination of the \(R \to 0\)
Limit}{4. Examination of the R \textbackslash to 0 Limit}}\label{examination-of-the-r-to-0-limit}

\subsubsection{\texorpdfstring{4.1 Divergence of
\(\Lambda\)}{4.1 Divergence of \textbackslash Lambda}}\label{divergence-of-lambda}

In the limit \(R \to 0\), \(\Lambda = 3/R^2 \to \infty\). The
\(r^2/R^2\) term in equation (3.4) becomes dominant, yielding

\[f(r) \approx -\frac{r^2}{R^2} \tag{4.1}\]

for any value of \(r > 0\). Since \(f(r) < 0\), surfaces of
\(r = \text{const}\) become timelike (corresponding to structures inside
a horizon).

\subsubsection{4.2 Horizon Structure}\label{horizon-structure}

The solutions of \(f(r) = 0\) give the positions of horizons. From
equation (3.4):

\[1 - \frac{2M}{r} - \frac{r^2}{R^2} = 0 \tag{4.2}\]

When \(R\) is sufficiently large, two positive solutions exist:

\begin{itemize}
\tightlist
\item
  \textbf{Black hole horizon} \(r_{\mathrm{BH}}\): near \(r \approx 2M\)
\item
  \textbf{Cosmological horizon} \(r_{\mathrm{C}}\): near \(r \approx R\)
\end{itemize}

In the limit \(R \to 0\), \(r_{\mathrm{C}} \to 0\), and the cosmological
horizon approaches and merges with the black hole horizon. That is, the
``region in which an observer can exist'' between the two horizons
vanishes.

\subsubsection{4.3 Correspondence with the Structural Tension of Paper
{[}4{]}}\label{correspondence-with-the-structural-tension-of-paper-4}

The structural tension identified in paper {[}4{]}, §3.2 --- where the
upper bound on geodesic deviation \(|\xi|^2_{\max}(R)\) approaches zero
while the local curvature diverges --- is expressed in the language of
the exact solution of this section as follows:

As \(R \to 0\), the cosmological horizon \(r_{\mathrm{C}} \to 0\), and
the region in which an internal observer can exist vanishes. This is the
exact-solution-level counterpart of the statement in paper {[}4{]}, §3.2
that ``the quantities measurable by an observer within the subjective
space approach zero.''

\subsubsection{4.4 Scope of the Claims in This
Section}\label{scope-of-the-claims-in-this-section}

The claims of this section are limited to the following:

\begin{enumerate}
\def\labelenumi{\arabic{enumi}.}
\tightlist
\item
  What geometric structure the exact solution (3.4) of equation (1.1)
  from paper {[}1{]} exhibits in the limit \(R \to 0\).
\item
  How this structure corresponds to the structural tension in paper
  {[}4{]}, §3.2.
\end{enumerate}

This paper does not claim that this exact solution describes any
specific astrophysical object (such as a black hole). \(M\) is a
mathematical integration constant, and its identification with physical
mass lies beyond the scope of this paper.

\begin{center}\rule{0.5\linewidth}{0.5pt}\end{center}

\subsection{5. Connection with the Nested
Structure}\label{connection-with-the-nested-structure}

\subsubsection{5.1 Application of Proposition
4.1}\label{application-of-proposition-4.1}

To the situation described in §4, where the region for an internal
observer vanishes in the \(R \to 0\) limit, Proposition 4.1 of paper
{[}4{]} (existence of nested structures) is applicable. For any value of
\(R > 0\), from any point within the subjective space \(\Pi_R\), a new
subjective space \(\Pi_{R'}\) with curvature radius \(R' > 0\) can be
constructed, where \(R'\) can be chosen independently of \(R\).

\subsubsection{5.2 Exact Solution in Child Subjective
Spaces}\label{exact-solution-in-child-subjective-spaces}

In the child subjective space \(\Pi_{R'}\), the construction of paper
{[}1{]} applies directly. Therefore, the vacuum Einstein equation in
\(\Pi_{R'}\) is

\[G'_{\mu\nu} + \Lambda'\, g'_{\mu\nu} = 0, \qquad \Lambda' = \frac{3}{R'^2} \tag{5.1}\]

and its spherically symmetric static exact solution is

\[f'(r') = 1 - \frac{2M'}{r'} - \frac{r'^2}{R'^2} \tag{5.2}\]

Since \(R'\) is independent of \(R\), even when \(R \to 0\) in the
parent subjective space, \(R'\) can be kept finite in the child
subjective space.

\subsubsection{5.3 Recovery of the Horizon
Structure}\label{recovery-of-the-horizon-structure}

The ``region in which an observer can exist,'' which vanished in the
parent subjective space as described in §4.2, is recovered within the
child subjective space \(\Pi_{R'}\) as long as \(R' > 0\) remains
finite. The exact solution (5.2) of the child subjective space possesses
a black hole horizon and a cosmological horizon according to the value
of \(R'\), and a well-defined region in which an observer can exist lies
between the two horizons.

\subsubsection{5.4 Scope of the Claims in This
Section}\label{scope-of-the-claims-in-this-section-1}

The claims of this section are limited to the following:

\begin{enumerate}
\def\labelenumi{\arabic{enumi}.}
\tightlist
\item
  By Proposition 4.1 of paper {[}4{]}, a child subjective space can be
  constructed from the domain of the \(R \to 0\) limit.
\item
  The equation from paper {[}1{]} and its exact solution hold also in
  the child subjective space.
\item
  The inter-horizon region that vanished in the parent subjective space
  is recovered within the child subjective space.
\end{enumerate}

These are direct consequences of Proposition 6.1 of paper {[}1{]} and
Proposition 4.1 of paper {[}4{]}, and contain no new assumptions.
Physical interpretation is outside the scope of this paper and is left
to the reader.

\begin{center}\rule{0.5\linewidth}{0.5pt}\end{center}

\subsection{6. Discussion}\label{discussion}

In this paper, we constructed the Schwarzschild--de Sitter metric as the
spherically symmetric static exact solution of the vacuum Einstein
equation derived in Proposition 6.1 of paper {[}1{]}, and connected the
behavior in the \(R \to 0\) limit with the nested structure of paper
{[}4{]}.

The descriptions in this paper are limited to the following:

\begin{enumerate}
\def\labelenumi{\arabic{enumi}.}
\tightlist
\item
  From equation (1.1) of paper {[}1{]} and Birkhoff's theorem, the
  Schwarzschild--de Sitter solution is uniquely determined without
  additional assumptions.
\item
  By substituting \(\Lambda = 3/R^2\), the exact solution is described
  in a concrete form with \(R\) as a parameter.
\item
  In the limit \(R \to 0\), the region for an internal observer
  vanishes, which constitutes the exact-solution-level counterpart of
  the structural tension in paper {[}4{]}, §3.2.
\item
  Through the nested structure of Proposition 4.1 of paper {[}4{]}, the
  inter-horizon region is recovered in the child subjective space.
\end{enumerate}

This paper does not claim that this exact solution describes any
specific physical situation. \(M\) is an integration constant and \(R\)
is the curvature radius of the central projection; the identification of
either with physical quantities lies beyond the scope of this paper.
Correspondences with physics concepts related to the Schwarzschild--de
Sitter metric (black holes, cosmological horizons, Hawking radiation,
information paradox, etc.) are problems outside the geometric framework
and are not addressed in this paper.

\subsubsection{Appendix: Suggested Logical Reading
Order}\label{appendix-suggested-logical-reading-order}

The papers in this series are numbered in order of publication, but the
following logical reading order is recommended:

\begin{itemize}
\tightlist
\item
  Papers {[}1{]}--{[}3{]}: Foundational framework
\item
  Paper {[}4{]}: Limits of \(R\) and nested structures
\item
  Paper {[}5{]}: Dimension addition
\item
  \textbf{This paper (Paper {[}9{]}): Exact solution and horizon
  structure in the \(R \to 0\) limit}
\item
  Paper {[}7{]}: Bridging of geodesics
\item
  Paper {[}6{]}: Description of subjective spaces in each dimension
\item
  Paper {[}8{]}: Dimensional development from zero dimensions
\end{itemize}

\begin{center}\rule{0.5\linewidth}{0.5pt}\end{center}

\subsection{7. Conclusion}\label{conclusion}

The spherically symmetric static exact solution of the vacuum Einstein
equation \(G_{\mu\nu} + \Lambda g_{\mu\nu} = 0\) (\(\Lambda = 3/R^2\))
in the central projection framework is uniquely determined by Birkhoff's
theorem (with \(\Lambda\) extension) as the Schwarzschild--de Sitter
metric \(f(r) = 1 - 2M/r - r^2/R^2\). No additional assumptions are
required for this construction. In the limit \(R \to 0\), the region for
an internal observer vanishes, but through the nested structure of
Proposition 4.1 of paper {[}4{]}, the same exact solution structure is
recovered in the child subjective space with \(R'\) as an independent
parameter. Physical interpretation is outside the scope of this paper
and is left to the reader.

\begin{center}\rule{0.5\linewidth}{0.5pt}\end{center}

\subsection{References}\label{references}

{[}1{]} Noriaki Kihara, ``Geometric Formalization of 4-Dimensional
Spacetime via Central Projection,'' Zenodo, DOI: 10.5281/zenodo.19427780
(2026).

{[}2{]} Noriaki Kihara, ``Geometric Symmetries of the Central
Projection: Mathematical Foundations of the Multi-Axis Model,'' Zenodo,
DOI: 10.5281/zenodo.19434932 (2026).

{[}3{]} Noriaki Kihara, ``Relativity of Observation in Multiple
Subjective Spaces: Geometric Consequences of Central Projection
Symmetries,'' Zenodo, DOI: 10.5281/zenodo.19435162 (2026).

{[}4{]} Noriaki Kihara, ``On the Limits of the Curvature Radius in
Central Projection --- The Cases \(R \to \infty\) and \(R \to 0\),''
Zenodo, DOI: 10.5281/zenodo.19526549 (2026).

{[}5{]} Noriaki Kihara, ``On Dimension Addition to Background Spaces and
Subjective Spaces,'' Zenodo, DOI: 10.5281/zenodo.19526913 (2026).

{[}6{]} Noriaki Kihara, ``Subjective Spaces from Zero to Four
Dimensions,'' Zenodo, DOI: 10.5281/zenodo.19533292 (2026).

{[}7{]} Noriaki Kihara, ``On the Continuity of Geodesics in Subjective
Spaces --- Bridging Geodesics from Parent to Child Subjective Spaces,''
Zenodo, DOI: 10.5281/zenodo.19533299 (2026).

{[}8{]} Noriaki Kihara, ``Diameter of the Circumscribed Hypersphere of a
Unit Four-Dimensional Hyperrectangle,'' Zenodo, DOI:
10.5281/zenodo.19533313 (2026).

\begin{center}\rule{0.5\linewidth}{0.5pt}\end{center}

\end{document}
